\scanpage{203}\mspage{152}
$\varphi_Q\circ i_2=f_{\unsure{2|\nu}}$, on voit déjà que

\[
\begin{aligned}
f_1\text{ et }f_2\text{ lisses}\Longleftrightarrow {}&\varphi_Q\text{ lisse sur l'ouvert réunion}\\
&\text{de }i_1(X_1)\text{ et }i_2(X_2),\ \text{i.e.}\\
&\text{sur }\unsure{\mathscr X|\mathcal X}\setminus i_0(X_0).
\end{aligned}
\]

Il reste à exprimer la lissité de $\varphi_Q$ en les points $x_0$ de $i_0(X_0)$. Soient $x_1,x_2,y$ les images de \deleted{\unsure{$x_0$ point}} dans $X_1,X_2,Y$ respectivement. Il faut exprimer que les fibres de \unsure{les} $(\varphi_Q)_{/x_0}\quad\mathscr X_{/x_0}\longrightarrow Y_{/y}$ sont asphériques. \deleted{Mais} Nous allons pour cela expliciter $\mathscr X_{/x_0}$ \inserted{et la flèche \unsure{can|sus}}, en termes du carré commutatif

\diagram{\((**)\qquad\begin{tikzcd}[ampersand replacement=\&,row sep=1.7em,column sep=3em]
X_{0/x_0}\arrow[r,"{\alpha_{1/x_0}}"]\arrow[d,"{\alpha_{2/x_0}}"'] \& X_{1/x_1}\arrow[d,"{f_{1/x_1}}"]\\
X_{2/x_2}\arrow[r,"{f_{2/x_2}}"] \& Y_{/y}
\end{tikzcd}\)}

induit par $Q$. Je dis

\uuline{Lemme 1} \unsure{Soit|Soient} \deleted{\unsure{$x_0$}}un carré commutatif \deleted{$Q$}

\diagram{\(\unsure{\in|\epsilon}\qquad Q=\deleted{\unsure{Q\to}}\begin{tikzcd}[ampersand replacement=\&,row sep=1em,column sep=2em]
X_0\arrow[r,"\alpha_1"]\arrow[d,"\alpha_2"'] \& X_1\arrow[d,dash,"f_1"]\\
X_2\arrow[r,"f_2"] \& Y
\end{tikzcd}\)}

dans $\mathrm{Cat}$, $x_0\in\mathrm{Ob}\,X_0$, $x_1,x_2,y$ ses images dans $X_1,X_2,Y$ respectivement, $Q_{/x_0}$ le carré \unsure{$(**)$|corresp.} ci-dessus, induit par $Q$ sur les localisés $X_{0/x_0}$ etc.,

\diagram{\(\mathscr X=\int \begin{tikzcd}[ampersand replacement=\&,row sep=0pt,column sep=1.2em,cramped] \& X_1 \\ X_0\arrow[ur]\arrow[dr,dash] \& \\ \& X_2 \end{tikzcd}.\)}

On a alors un
% archive sheet number: 203
